% P11 paper module: complete proof of the sharp fixed-vector odd Schur asymptotic.
% Canonical inputs condensed here: C4, O3c, C5c, O3d-I1, O3d-I2.

\begin{proof}[Proof of Theorem~\ref{thm:odd}]
Put
\[
 h_T:=H_T^*J_{R,T}f_-,
 \qquad
 A_T:=I+R_T^*R_T.
\]
By parity, $h_T$ is even.  The full-rest factorization
\eqref{eq:full-rest-gram}, now at terminal level $T$, gives a canonical analysis
operator $\widetilde R_T:\mathscr H_T\to\mathscr Z_T$ with
\[
 A_T=I+\widetilde R_T^*\widetilde R_T.
\]
Completing the square yields the exact dual formula
\[
 \boxed{
 \langle h_T,A_T^{-1}h_T\rangle
 =\inf_{Y\in\mathscr Z_T}
 \bigl(\|h_T-\widetilde R_T^*Y\|_2^2+\|Y\|_{\mathscr Z_T}^2\bigr).
 }
\tag{6.2}\label{eq:odd-dual}
\]
Indeed, the minimizer is
$Y=(I+\widetilde R_T\widetilde R_T^*)^{-1}\widetilde R_Th_T$.

\paragraph{Step 1: matching lower scale.}
By Theorem~\ref{thm:jet-expansion}, if $m=m(f_-)$ is the first nonzero integral jet,
\[
 \ell_T(f_-):=\langle h_T,\mathbf1_T\rangle
 =-\sqrt2\,c_m\beta_R^{(m)}(f_-)
 \frac{e^{T/2}}{T^{m+1/2}}\bigl(1+O(T^{-1})\bigr).
\tag{6.3}\label{eq:odd-ell-leading}
\]
The sharp constant-mode estimate is
\[
 \boxed{
 d_T:=\langle\mathbf1_T,A_T\mathbf1_T\rangle=2T+O(1).
 }
\tag{6.4}\label{eq:odd-constant-sharp}
\]
For completeness, the primitive conditioned rest vanishes on $\mathbf1_T$, while the
higher prime-power part is absolutely summable in $p,k$; its total $L^2$ contribution
is $O(1)$.  Thus
$\|R_T\mathbf1_T\|^2=O(1)$ and~\eqref{eq:odd-constant-sharp} follows from
$\|\mathbf1_T\|_2^2=2T$.
The variational inequality
\[
 \langle h,A_T^{-1}h\rangle
 \ge\frac{|\langle h,e\rangle|^2}{\langle e,A_Te\rangle}
\]
with $e=\mathbf1_T$ gives
\[
 \sigma_T(J_{R,T}f_-)
 \ge c_m^2|\beta_R^{(m)}(f_-)|^2
 \frac{e^T}{T^{2m+2}}\bigl(1+O(T^{-1})\bigr).
\tag{6.5}\label{eq:odd-lower-sharp}
\]
It remains to construct a dual certificate with the same leading cost.

\paragraph{Step 2: growing primitive hub and bounded remainder.}
Choose $\rho_f<R$ with
$\operatorname{supp}f_-\subset(-\rho_f,\rho_f)$ and fix
$a_*>\rho_f+2\varepsilon$, where
$0<\varepsilon<\frac12\log2$.  For a primitive prime $p$ write
\[
 a_p:=\frac12\log p,
 \qquad
 d_p:=T-a_p,
 \qquad
 c_p:=\sqrt{\log p}\,p^{-3/4}.
\]
Set
\[
 h_T^{\rm grow}
 :=\sum_{p:\,a_*\le a_p\le T}
 c_pD_{\log p}^*J_{R,T}f_-,
 \qquad
 h_T^{\rm rem}:=h_T-h_T^{\rm grow}.
\tag{6.6}\label{eq:odd-hub-split}
\]
The first omitted primitive block contains only finitely many primes.  For all higher
prime powers, $\|D_s\|\le2$ and
\[
 \sum_p\sum_{k\ge2}\sqrt{\log p}\,p^{-3k/4}<\infty.
\]
Hence
\[
 \boxed{\sup_T\|h_T^{\rm rem}\|_2<\infty,}
 \qquad
 |\langle h_T^{\rm rem},\mathbf1_T\rangle|=O(\sqrt T).
\tag{6.7}\label{eq:odd-rem-bounded}
\]

For an even $e\in\mathscr H_T$ put $b_T(t):=e(T-t)$, $0<t<T$.  A direct change of
variables gives
\[
 \langle h_T^{\rm grow},e\rangle
 =\int_0^Tk_T(t)b_T(t)\,dt,
\tag{6.8}\label{eq:odd-kernel-pairing}
\]
where
\[
 k_T(t):=-2\sum_{p:\,a_*\le a_p\le T}c_p f_-(t-d_p).
\tag{6.9}\label{eq:odd-kernel}
\]
For $A=T-t$, ordinary partial summation using the prime number theorem and the compact
support of $f_-$ and $f_-'$ gives uniformly
\[
 |k_T(t)|+|k_T'(t)|
 \le C_{R,f_-}\frac{e^{A/2}}{\sqrt{1+A}}.
\tag{6.10}\label{eq:odd-kernel-bound}
\]
Consequently
\[
 e^{-T}\int_0^T e^{t/2}|k_T(t)|^2dt
 =O(T^{-1})+O(e^{-T/4}).
\tag{6.11}\label{eq:odd-kernel-weighted}
\]

\paragraph{Step 3: exact mean-zero splitting.}
Let
\[
 K_T:=\int_0^Tk_T(t)dt
 =\langle h_T^{\rm grow},\mathbf1_T\rangle.
\]
Equations~\eqref{eq:odd-ell-leading} and~\eqref{eq:odd-rem-bounded} imply
\[
 K_T=-\sqrt2\,c_m\beta_R^{(m)}(f_-)
 \frac{e^{T/2}}{T^{m+1/2}}(1+o(1)).
\tag{6.12}\label{eq:odd-KT}
\]
Define
\[
 \mu_T:=\frac{K_T}{2T},
 \qquad
 g_T^{\rm grow}:=h_T^{\rm grow}-\mu_T\mathbf1_T,
 \qquad
 k_T^0(t):=k_T(t)-\frac{K_T}{T}.
\tag{6.13}\label{eq:odd-mean-zero}
\]
Then
\[
 \langle g_T^{\rm grow},\mathbf1_T\rangle=0,
 \qquad
 \int_0^Tk_T^0(t)dt=0.
\]
The cost of the extracted constant mode is
\[
 M_T:=\|\mu_T\mathbf1_T\|_2^2
 =\frac{|K_T|^2}{2T}
 =c_m^2|\beta_R^{(m)}(f_-)|^2
 \frac{e^T}{T^{2m+2}}(1+o(1)).
\tag{6.14}\label{eq:odd-MT}
\]

\paragraph{Step 4: signed continuous future-edge certificate.}
Choose
$\alpha\in C_c^\infty((0,\varepsilon))$ with $\int\alpha=1$.
Since $k_T^0$ has mean zero,
\[
 \int_0^Tk_T^0(t)b(t)dt
 =\int_0^T\int_0^\varepsilon
 k_T^0(t)\alpha(s)[b(t)-b(s)]\,ds\,dt.
\tag{6.15}\label{eq:odd-anchor-identity}
\]
With $r=(t+s)/2$ this becomes
\[
 \int k_T^0b
 =\iint C_T^-(r,t)[b(t)-b(2r-t)]\,dt\,dr,
 \qquad
 C_T^-(r,t):=2k_T^0(t)\alpha(2r-t).
\tag{6.16}\label{eq:odd-signed-continuous}
\]
For a primitive prime $q$, with $r_q=T-\frac12\log q$, the corresponding conditioned
primitive edge on even terminal vectors is exactly
\[
 b(t)-b(2r_q-t).
\]
Thus~\eqref{eq:odd-signed-continuous} is the continuous model of the actual primitive
future-rest adjoint.

The continuous primitive rest mass density is
$m_T(r)=2e^{T-r}$.  Dividing $C_T^-$ by $m_T(r)^{1/2}$ and changing variables
$s=2r-t$ gives
\[
 \|Y_T^{\rm cont,-}\|^2
 \le C_\alpha e^{-T}\int_0^Te^{t/2}|k_T^0(t)|^2dt.
\]
Using~\eqref{eq:odd-kernel-weighted} and
$|k_T^0|^2\le2|k_T|^2+2|K_T|^2/T^2$ yields
\[
 \boxed{\|Y_T^{\rm cont,-}\|^2=o(M_T).}
\tag{6.17}\label{eq:odd-cont-cost}
\]

\paragraph{Step 5: discrete prime quadrature.}
We use the unconditional prime number theorem in short intervals in the following
fixed form: for every fixed $\theta>17/30$,
\[
 \sum_{x<n\le x+x^\theta}\Lambda(n)\sim x^\theta
 \qquad(x\to\infty).
\tag{6.18}\label{eq:short-PNT}
\]
This is the Guth--Maynard short-interval consequence; we take $\theta=3/5$.
At prime scale $q\asymp X=e^{2(T-r)}$, such intervals correspond to $r$-cells $I$ of
width
\[
 |I|\asymp X^{-2/5}=e^{-\frac45(T-r)}.
\]
For
\[
 w_q:=\frac{\log q}{\sqrt q}\left(1-\frac1q\right),
 \qquad
 W_I:=\sum_{q:r_q\in I}w_q,
\]
short-interval PNT gives uniformly on the future cells
\[
 \boxed{W_I\asymp m_T(r_I)|I|.}
\tag{6.19}\label{eq:prime-cell-mass}
\]
Define the mass-normalized coefficients
\[
 \lambda_q^{(I)}:=|I|\frac{w_q}{W_I}.
\]
Then exactly
\[
 \sum_{q:r_q\in I}\lambda_q^{(I)}=|I|,
 \qquad
 \sum_{q:r_q\in I}\frac{|\lambda_q^{(I)}|^2}{w_q}
 =\frac{|I|^2}{W_I}
 \lesssim\frac{|I|}{m_T(r_I)}.
\tag{6.20}\label{eq:prime-cell-cost}
\]
For every Hilbert-valued Lipschitz function $\Phi$,
\[
 \left\|\sum_{q:r_q\in I}\lambda_q^{(I)}\Phi(r_q)
 -\int_I\Phi(r)dr\right\|
 \le2|I|^2\sup_{r\in I}\|\Phi'(r)\|.
\tag{6.21}\label{eq:prime-cell-quadrature}
\]
Apply this to the signed edge field in~\eqref{eq:odd-signed-continuous}.  The cost
identity~\eqref{eq:prime-cell-cost} and~\eqref{eq:odd-cont-cost} produce a primitive
discrete certificate $Y_T^{\rm prim,-}$ with
\[
 \boxed{\|Y_T^{\rm prim,-}\|^2=o(M_T).}
\tag{6.22}\label{eq:odd-prim-cost}
\]
The derivative bound~\eqref{eq:odd-kernel-bound} and
$\max_I|I|\ll e^{-2T/5}$ give the source quadrature remainder
\[
 \boxed{\|Z_T^{\rm quad}\|_2=o(\sqrt{M_T}).}
\tag{6.23}\label{eq:odd-quadrature-error}
\]
Indeed,~\eqref{eq:prime-cell-quadrature} summed over the $O(T)$ total $r$-length gives an
exponential factor at most $e^{-cT}$ for the $k_T$ part; the constant part
$-K_T/T$ has zero $t$-derivative and contributes
$O(|K_T|e^{-2T/5})=o(\sqrt{M_T})$.

\paragraph{Step 6: full-rest lift; no primitive form domination.}
The primitive certificate of Step~5 is not inserted into the full Schur problem by an
operator inequality.  Instead use the exact $a=0$ channel of $\widetilde R_T$.
For each future prime $q$ it decomposes as
\[
 (\widetilde R_Tf)_{q,0}
 =\underbrace{\sqrt{w_q}\,1_{\Omega_{q,0,T}}
 K_{\log q}^{\rm tr}f}_{P_{q,T}f}
 +\underbrace{\sqrt{(\log q)(q-1)}\,1_{\Omega_{q,0,T}}
 \sum_{k\ge2}q^{-3k/4}K_{k\log q}^{\rm tr}f}_{E_{q,T}f}.
\tag{6.24}\label{eq:future-primitive-tail}
\]
Here $K_s^{\rm tr}=P_TD_sE_T$.  Since $\|K_s^{\rm tr}\|\le2$,
\[
 \|E_{q,T}\|\le C\sqrt{\log q}\,q^{-1}.
\]
All cells used above satisfy
$\frac12\log q\ge T/2-O(1)$, hence
\[
 \boxed{
 \|E_T^{\rm fut}\|
 \le C\sqrt{T+1}\,e^{-T/2}.
 }
\tag{6.25}\label{eq:future-tail-small}
\]
Lift $Y_T^{\rm prim,-}$ isometrically into the $a=0$ component of $\mathscr Z_T$ and
call the lift $\widehat Y_T^-$.  Absorbing the tail
$(E_T^{\rm fut})^*\widehat Y_T^-$ into the source remainder gives
\[
 g_T^{\rm grow}
 =\widetilde R_T^*\widehat Y_T^-
 +Z_T^{\rm quad}+Z_T^{\rm tail},
\tag{6.26}\label{eq:full-rest-lift}
\]
with
\[
 \|\widehat Y_T^-\|^2=o(M_T),
 \qquad
 \|Z_T^{\rm tail}\|_2=o(\sqrt{M_T}).
\tag{6.27}\label{eq:full-rest-lift-cost}
\]
This is the O3d-I1 repair in the form needed here: the unproved order
$R_T^*R_T\ge(R_T^{(1)})^*R_T^{(1)}$ is never used.

\paragraph{Step 7: full-hub squeeze.}
From~\eqref{eq:odd-hub-split},~\eqref{eq:odd-mean-zero} and
\eqref{eq:full-rest-lift},
\[
 h_T=\widetilde R_T^*\widehat Y_T^-+Z_T^{\rm full},
\]
where
\[
 Z_T^{\rm full}
 :=\mu_T\mathbf1_T+h_T^{\rm rem}+Z_T^{\rm quad}+Z_T^{\rm tail}.
\]
Equations~\eqref{eq:odd-rem-bounded},~\eqref{eq:odd-quadrature-error} and
\eqref{eq:full-rest-lift-cost}, together with $\sqrt{M_T}\to\infty$, imply
\[
 \|Z_T^{\rm full}\|_2^2=M_T(1+o(1)),
 \qquad
 \|\widehat Y_T^-\|^2=o(M_T).
\tag{6.28}\label{eq:full-hub-cost}
\]
Insert this admissible pair into the dual formula~\eqref{eq:odd-dual}:
\[
 \sigma_T(J_{R,T}f_-)
 \le M_T(1+o(1)).
\tag{6.29}\label{eq:odd-upper-sharp}
\]
Together with~\eqref{eq:odd-lower-sharp} and~\eqref{eq:odd-MT}, this gives exactly
\eqref{eq:odd-asymptotic}.
\end{proof}

\begin{remark}[Scope of the sharp asymptotic]\label{rem:odd-fixed-vector-firewall}
The proof is for each fixed smooth odd source vector.  The $o(1)$ is not uniform on the
odd unit sphere and gives no direct control of $G_{R,T}^{-1/2}$ on $T$-dependent
vectors.  In particular, it does not prove strong terminal transport.
\end{remark}
